\subsection{Height one quartics}
\begin{center}\small
\begin{tabular}{@{}lrrrr@{}}\toprule
Endpoint & Variables & Terms & Height & Constant\\\midrule
\file{compact571} &571&19833&1&1\\
\file{balanced593} &593&12027&1&1\\
\file{sparse646} &646&8368&1&1\\
\file{sparser742} &742&7226&1&1\\\bottomrule
\end{tabular}
\end{center}
All variables are unrestricted integers. The first two graphs use the new
263-variable signed source; the last two use the 285-variable source.
Every collected nonzero coefficient is $+1$ or $-1$ and every declared
variable occurs. The 571-variable point improves both dimension and support
over the Round~17 793- and 806-variable height-one points; the
646-variable point does the same relative to its 994-variable sparse point.
The four new choices retain different dimension/support tradeoffs.
These are coefficient-height results, not claims of smaller dimension than
the direct signed quartics.
